% ============================================================
% SECȚIUNEA 4 — Pipeline-ul de feedback AI
% Fișiere sursă de referință:
%   backend/AIInterviewCoach.Application/Services/SubmissionFeedbackProcessor.cs
%   backend/AIInterviewCoach.Infrastructure/Services/SubmissionFeedbackBackgroundService.cs
%   backend/AIInterviewCoach.Infrastructure/Services/SubmissionFeedbackService.cs
%   frontend/src/app/features/candidate/services/candidate-workspace.facade.ts
%   backend/AIInterviewCoach.Domain/Constants/SubmissionFeedbackStatuses.cs
% ============================================================

\section{Pipeline-ul de feedback AI}
\label{sec:pipeline_feedback_ai}

\subsection{Vedere de ansamblu}

Evaluarea cu ajutorul inteligenței artificiale funcționează \emph{asincron}
față de cererea HTTP care creează submisia. Decizia de a nu bloca
răspunsul HTTP (care ar putea dura câteva secunde) se justifică prin
latența imprevizibilă a API-ului OpenAI și prin nevoia de a proteja
experiența candidatului. Submisia este creată imediat cu statusul
\texttt{Pending}; feedback-ul sosește ulterior, iar frontend-ul îl afișează
prin polling.

Pipeline-ul este alcătuit din patru componente principale, descrise în
secțiunile care urmează:

\begin{enumerate}
  \item \texttt{SubmissionService} — creează submisia și o pune în coadă.
  \item \texttt{SubmissionFeedbackBackgroundService} — coadă de fundal
        bazată pe \texttt{System.Threading.Channels}.
  \item \texttt{SubmissionFeedbackProcessor} — orchestrează generarea și
        persistă rezultatul.
  \item \texttt{SubmissionFeedbackService} — apelează efectiv API-ul
        OpenAI sau furnizează feedback local de rezervă.
\end{enumerate}

\subsection{Mașina de stări a feedback-ului}

Coloana \texttt{AiFeedbackStatus} din tabela \texttt{Submissions} poate
lua trei valori, definite ca constante de șir în
\texttt{SubmissionFeedbackStatuses}:
% Sursă: Domain/Constants/SubmissionFeedbackStatuses.cs

\begin{itemize}
  \item \texttt{Pending} — stare inițială (valoare implicită la inserție
        în baza de date); feedback-ul nu a fost generat încă.
  \item \texttt{Ready} — feedback-ul a fost generat cu succes și este
        stocat în coloana \texttt{AiFeedback}.
  \item \texttt{Failed} — generarea a eșuat definitiv (eroare API
        neașteptată sau problemă lipsă din baza de date).
\end{itemize}

Tranziția \texttt{Pending}$\rightarrow$\texttt{Ready} sau \texttt{Pending}$\rightarrow$\texttt{Failed}
este atomică: ambele câmpuri (\texttt{AiFeedback} și \texttt{AiFeedbackStatus})
sunt actualizate într-un singur apel \texttt{SaveChangesAsync()}.
% Sursă: Application/Services/SubmissionFeedbackProcessor.cs:75–86

\subsection{SubmissionFeedbackBackgroundService — coada de procesare}

\texttt{SubmissionFeedbackBackgroundService} este un
\texttt{IHostedService} (singleton) care implementează totodată
\texttt{ISubmissionFeedbackQueue}, interfața prin care serviciile de
aplicație înregistrează submisii în coadă.
% Sursă: Infrastructure/Services/SubmissionFeedbackBackgroundService.cs

\paragraph{Mecanismul de coadă}
Se folosesc două structuri de date paralele:
\begin{itemize}
  \item \texttt{Channel<Guid>} — canal neblocat, configurat cu
        \texttt{SingleReader = true, SingleWriter = false}; vehiculează
        ID-urile submisiilor de procesat.
  \item \texttt{ConcurrentDictionary<Guid, byte>} — set de deduplicare;
        previne adăugarea aceleiași submisii de două ori (relevant dacă
        un ID ajunge din nou în coadă înainte de a fi procesat).
\end{itemize}

Metoda \texttt{QueueAsync()} verifică mai întâi dicționarul; dacă ID-ul
este deja prezent, returnează imediat fără a scrie în canal.
% Sursă: Infrastructure/Services/SubmissionFeedbackBackgroundService.cs:34–42

\paragraph{Recuperare la pornire}
La pornirea aplicației, \texttt{ExecuteAsync()} interoghează baza de date
și pune în coadă toate submisiile cu \texttt{AiFeedbackStatus = ``Pending''},
ordonate după \texttt{SubmittedAt}. Această recuperare rezolvă cazul
în care serverul este repornit cu submisii nerezolvate.
% Sursă: Infrastructure/Services/SubmissionFeedbackBackgroundService.cs:56–81

\paragraph{Bucla de procesare}
Bucla principală citește ID-uri din canal cu
\texttt{WaitToReadAsync()} (blocare asincronă fără busy-wait) și,
pentru fiecare ID, creează un scope DI nou, instanțiează
\texttt{SubmissionFeedbackProcessor} și îl apelează. Excepțiile
neașteptate sunt loggate (nivel \texttt{Error}) dar nu opresc bucla;
\texttt{OperationCanceledException} la oprire iese curat.
% Sursă: Infrastructure/Services/SubmissionFeedbackBackgroundService.cs:44–77

\subsection{SubmissionFeedbackProcessor — orchestrare}

\texttt{SubmissionFeedbackProcessor} este un serviciu \emph{scoped} din
stratul Application; primește un singur ID de submisie și execută
urmăroarele etape în ordine:
% Sursă: Application/Services/SubmissionFeedbackProcessor.cs

\begin{enumerate}
  \item \textbf{Interogare}: Încarcă submisia împreună cu entitatea
        \texttt{Problem} asociată printr-un \texttt{Include(s => s.Problem)}.
        Dacă submisia nu mai există (ștearsă între timp), abandonează
        silențios.
        % Sursă: SubmissionFeedbackProcessor.cs:33–41

  \item \textbf{Verificare stare}: Dacă \texttt{AiFeedbackStatus} nu mai
        este \texttt{Pending} (submisie deja procesată), returnează imediat
        fără efect.
        % Sursă: SubmissionFeedbackProcessor.cs:43–47

  \item \textbf{Apel serviciu}: Apelează
        \texttt{ISubmissionFeedbackService.GenerateFeedbackAsync()} cu un
        DTO de context ce include titlul problemei, codul sursă, rezultatul
        execuției, limbajul și modul de utilizare (practică vs.\ interviu).
        % Sursă: SubmissionFeedbackProcessor.cs:56–73, BuildFeedbackContext()

  \item \textbf{Persistare}: Actualizează \texttt{AiFeedback} și
        \texttt{AiFeedbackStatus} și apelează \texttt{SaveChangesAsync()}.
        % Sursă: SubmissionFeedbackProcessor.cs:75–86

  \item \textbf{Observabilitate}: Înregistrează metrica de latență și
        sursa feedback-ului (\texttt{OpenAI} sau \texttt{LocalFallback})
        prin \texttt{IObservabilityService.RecordAiFeedback()}.
        % Sursă: SubmissionFeedbackProcessor.cs:88–92
\end{enumerate}

\subsection{SubmissionFeedbackService — generarea efectivă}

\texttt{SubmissionFeedbackService} este responsabil de comunicarea cu
API-ul OpenAI și de logica de rezervă.
% Sursă: Infrastructure/Services/SubmissionFeedbackService.cs

\paragraph{Configurare}
Cheia API este citită prioritizat din:
\begin{enumerate}
  \item Cheia de configurare \texttt{AiFeedback:ApiKey}
        (\texttt{appsettings.json} / variabile de mediu).
  \item Variabila de mediu \texttt{OPENAI\_API\_KEY}.
\end{enumerate}
Modelul este obținut similar din \texttt{AiFeedback:Model} sau
\texttt{OPENAI\_MODEL}, iar URL-ul de bază din \texttt{AiFeedback:BaseUrl}
sau \texttt{OPENAI\_BASE\_URL} (implicit
\texttt{https://api.openai.com/v1}).
% Sursă: SubmissionFeedbackService.cs:78–101

Dacă oricare dintre cheie și model lipsesc, serviciul nu apelează API-ul
și returnează direct feedback-ul de rezervă.
% Sursă: SubmissionFeedbackService.cs:57–59

\paragraph{Apelul API}
Serviciul trimite un \texttt{POST} la \texttt{\{baseUrl\}/responses}
cu corp JSON care include:
\begin{itemize}
  \item \texttt{model} — identificatorul modelului configurat.
  \item \texttt{instructions} — prompt de sistem cu instrucțiunile de stil
        (text simplu, patru secțiuni: \emph{Overall}, \emph{Correctness},
        \emph{Code Quality}, \emph{Next Step}; fără Markdown).
  \item \texttt{input} — promptul candidat construit din contextul submisiei
        (titlu, descriere, constrângeri, cod sursă, output de execuție,
        număr de teste trecute, mod de utilizare).
  \item \texttt{max\_output\_tokens = 220} — limitare explicită pentru
        menținerea răspunsului concis.
\end{itemize}
% Sursă: SubmissionFeedbackService.cs:73–88, BuildPrompt()

\paragraph{Parsarea răspunsului}
Funcția \texttt{ExtractResponseText()} suportă două formate de răspuns
OpenAI: câmpul de nivel superior \texttt{output\_text} (răspuns simplificat)
și structura \texttt{output[*].content[*].text} (răspuns complet).
% Sursă: SubmissionFeedbackService.cs:106–135

\paragraph{Normalizare}
Funcția \texttt{NormalizeRemoteFeedback()} elimină heading-urile Markdown
(\texttt{\#}), elimină liniile goale de la început/sfârșit și colapsează
secvențele de mai multe linii goale consecutive la una singură.
Rezultatul este text simplu, gata de afișare în cardul UI.
% Sursă: SubmissionFeedbackService.cs:137–167

\paragraph{Feedback de rezervă}
În caz de eroare de rețea, de configurare lipsă sau de răspuns gol,
\texttt{SubmissionFeedbackService} returnează conținut generat de
\texttt{SubmissionFallbackFeedbackCatalog.Generate()}, un catalog de
texte statice personalizate după statusul submisiei.
Sursa este marcată ca \texttt{LocalFallback} în DTO, astfel încât
\texttt{SubmissionFeedbackProcessor} poate distinge cele două cazuri
pentru observabilitate.
% Sursă: SubmissionFeedbackService.cs:196–202

\subsection{Polling-ul din frontend}

Frontend-ul (Angular) detectează submisiile cu \texttt{aiFeedbackStatus = ``Pending''}
și programează o verificare periodică până când toate devin \texttt{Ready}
sau \texttt{Failed}.
% Sursă: frontend/.../candidate-workspace.facade.ts:912–1000

Mecanismul este implementat în \texttt{CandidateWorkspaceFacade}:

\begin{enumerate}
  \item La fiecare actualizare a listei de submisii, metoda
        \texttt{setSubmissions()} verifică dacă există cel puțin una cu
        \texttt{aiFeedbackStatus === ``Pending''}.
  \item Dacă da, apelează \texttt{ensureFeedbackPolling(delayMs = 1500)}.
        Această metodă creează un \texttt{setTimeout} de 1\,500\,ms dacă
        nu există deja unul activ (idempotentă).
  \item La expirarea timeout-ului, \texttt{refreshPendingFeedback()} face
        un apel API:
        \begin{itemize}
          \item \textbf{Mod practică}: \texttt{GET /api/submissions/me},
                filtrând local submisiile care corespund problemei curente.
          \item \textbf{Mod interviu}: \texttt{GET /api/submissions/session/\{sessionId\}}.
        \end{itemize}
  \item Răspunsul este pasat din nou prin \texttt{setSubmissions()}, care
        decide dacă polling-ul continuă sau se oprește.
  \item În caz de eroare a cererii, \texttt{ensureFeedbackPolling(4000)}
        reprogramează cu o întârziere mai mare de 4\,000\,ms.
  \item Polling-ul se oprește definitiv (\texttt{stopFeedbackPolling()})
        când nu mai există submisii \texttt{Pending} sau la distrugerea
        componentei, prin \texttt{DestroyRef}.
\end{enumerate}
% Sursă: candidate-workspace.facade.ts:929–999

Alegerea \texttt{setTimeout} în locul unui interval fix evită
acumularea cererilor dacă una este lentă: o nouă verificare nu
este programată decât după ce răspunsul anterior a sosit.
